% drafts/04-emergence.tex — Emergent Properties from Tree Structure
%
% Target length: ~5 pages (two-column PRA)
% Subsections: 4.1 Emergent locality, 4.2 Symmetry fractionalization,
%              4.3 Gauge structure, 4.4 Renormalization

This section is the conceptual heart of the paper.  We show that
locality, symmetry, gauge structure, and renormalization are not
intrinsic to the physics — they emerge from the tree.

%%====================================================================
\subsection{Emergent locality}
\label{sec:em-locality}

\begin{definition}[Encoding weight]
\label{def:encoding-weight}
For an encoding $E$ and a fermionic operator $O$, the
\emph{encoding weight} is
$\pauliw{E(O)} = \max_{\sigma \in \mathrm{supp}(E(O))} \pauliw{\sigma}$.
\end{definition}

\begin{theorem}[Weight bound]
\label{thm:weight-bound}
For a tree $T$ of depth~$d$ with path-based encoding,
\begin{equation}
  \pauliw{\ad{j}} \le 2\,\depth{j} + 1.
  \label{eq:weight-bound}
\end{equation}
\end{theorem}

\begin{proof}
In Construction~B, the Majorana string for a leg at the end of a
follow-$X/Y$-then-$Z$ path from node~$j$ has at most $\depth{j} + d_Z$
non-identity Pauli factors, where $d_Z$ counts the $Z$-links followed
after the initial $X$ or $Y$ step.  In the worst case
$d_Z \le \depth{j}$, and a ladder operator involves two Majorana
strings, giving $\pauliw{\ad{j}} \le 2\,\depth{j} + 1$.

See \cref{app:weight-proof} for the detailed argument.
\end{proof}

\begin{corollary}
For a balanced ternary tree with branching factor~3,
$\depth{j} \le \lfloor \log_3 n \rfloor$, so
$\pauliw{\ad{j}} = O(\log_3 n)$.
\end{corollary}

\paragraph{The locality paradox.}
Consider the nearest-neighbour hopping term $\ad{0}\an{1}$.
Under the Jordan--Wigner encoding this has weight~2 (two-local), but
under the Parity encoding it has weight~$n$ (maximally non-local).
The physics is identical — same eigenspectrum, same ground-state energy
— yet the \emph{manifest locality} is entirely different.

\begin{quote}\itshape
Locality is not a property of the interaction.\\
It is a property of the \emph{representation} of the interaction.
\end{quote}

This is precisely the sense in which locality is \emph{emergent}: it
does not exist in the algebraic (fermionic) description and is created
by the encoding.  The parallel to holographic reconstruction —  where a
bulk local operator can be a simple or complex boundary operator
depending on the code~\cite{pastawski2015happy} — is exact.

% Figure: same hopping term, weight comparison across 5 encodings.
% At n=8: JW weight-2, BK weight-4, Parity weight-8,
%         BinTree weight-4, TerTree weight-3.

%%====================================================================
\subsection{Emergent symmetry fractionalization}
\label{sec:em-symmetry}

The particle-number parity operator is a global $\mathbb{Z}_2$ symmetry:
%
\begin{equation}
  P = \prod_{j=0}^{n-1} (I - 2\nhat{j}) = \prod_j (-1)^{\nhat{j}}.
  \label{eq:parity-op}
\end{equation}

\begin{theorem}[Parity weight]
\label{thm:parity-weight}
Under the path-based encoding for a tree $T$, the encoded parity
operator $\hat{P} = E(P)$ satisfies
\begin{equation}
  \hat{P} = \prod_{v \in Z\text{-chain}(T)} Z_v,
\end{equation}
where $Z$-chain$(T)$ is the maximal path from the root following
$Z$-links.  Hence $\pauliw{\hat{P}} = |Z\text{-chain}(T)|$.
\end{theorem}

\paragraph{Examples.}
\begin{itemize}
  \item \textbf{JW} (star, linear ordering): $Z$-chain = full chain,
    weight~$n$.  The parity operator is $Z^{\otimes n}$.
  \item \textbf{Parity} (star, reverse ordering): $Z$-chain =
    root only, weight~1.  A single $Z_{n-1}$ encodes global parity.
  \item \textbf{BK} (Fenwick, $n = 2^k$): weight~1.  The root qubit
    carries total parity as~$Z_{n-1}$.
  \item \textbf{BinTree} (balanced binary): weight~1.  The root
    qubit carries parity as~$Z_{\text{root}}$.
  \item \textbf{TerTree} (balanced ternary): $Z$-chain =
    $\lfloor \log_3 n \rfloor + 1$ qubits.  For $n = 8$, weight~2;
    for $n = 27$, weight~3.
\end{itemize}

The same $\mathbb{Z}_2$ symmetry is represented as a weight-$n$ (global)
or weight-1 (local) Pauli operator, depending solely on the tree.  This
is a concrete instance of \emph{symmetry fractionalization}
\cite{wen2017zoo}: the tree determines how a global symmetry is
distributed among the qubits.

% Figure: parity operator across 5 encodings at n=8.
% JW: weight 8 (ZZZZZZZZ).  BK: weight 1 (Z_7).  Parity: weight 1 (Z_7).
% BinTree: weight 1 (Z_root).  TerTree: weight 2 (Z_root Z_child).

%%====================================================================
\subsection{Emergent gauge structure}
\label{sec:em-gauge}

When the physical Hilbert space is restricted — e.g., to states with
fixed particle number~$N_e$ — the encoding introduces stabiliser
constraints that act as emergent gauge symmetries.

\paragraph{Example ($n=4$, $N_e=2$).}
Under the Parity encoding (a star tree in reverse ordering),
$Z_3 = (-1)^{N_e}$.  For $N_e = 2$ (even),
$Z_3 = +1$, so qubit~3 is frozen: it can be eliminated by qubit
tapering~\cite{bravyi2017tapering}.  Under Jordan--Wigner (also a star,
with direct ordering), no single
qubit carries the particle-number parity, and no direct tapering is
available from this symmetry alone.

Different trees produce different stabiliser groups and hence different
effective theories after gauge fixing.  The tree determines not just
\emph{which} gauge constraints exist, but \emph{how many} qubits can
be eliminated.

% Stabiliser analysis: at n=8, Parity/BK/BinTree allow tapping 1 qubit
% via the weight-1 parity operator.  JW requires all 8 qubits.
% TerTree allows tapping via a weight-2 operator.

%%====================================================================
\subsection{The tree as renormalization scheme}
\label{sec:em-renorm}

In a balanced ternary tree:
\begin{itemize}
  \item Leaf nodes encode individual occupation numbers (microscopic).
  \item Internal nodes at depth~$d$ encode parity of $O(3^d)$ modes
        (mesoscopic).
  \item The root encodes total parity (macroscopic).
\end{itemize}
This hierarchy is a coarse-graining: each level aggregates three
sublevels — structurally identical to the multi-scale entanglement
renormalization ansatz (MERA)~\cite{vidal2007mera,swingle2012}.

\begin{conjecture}[MERA connection]
\label{conj:mera}
For a 1D fermionic system with a critical ground state, the tree
minimising total Pauli weight is structurally similar to the optimal
MERA network for the same system.
\end{conjecture}

\begin{conjecture}[Hamiltonian-adapted trees]
\label{conj:hamiltonian-adapted}
For a Hamiltonian $H$, the tree minimising
$\sum_{\text{terms}} |c_t| \cdot \pauliw{t}$ corresponds to a
renormalization scheme adapted to the entanglement structure of $H$'s
ground state.
\end{conjecture}

These conjectures are presented as motivation for future work.
The precise connection between encoding optimality and renormalization
optimality remains open.
